\section{Filesystem State Space and Invariants}\label{sec:state}

We now define formally the object on which the recovery operator acts.
The aim is to make every subsequent proof refer to a precisely typed
object rather than to an informal notion of ``filesystem state''.

\subsection{State Space}

\begin{definition}[Layered state]\label{def:state}
The BEAMFS state space~\(\mathcal{S}\) is the Cartesian product
\[
  \mathcal{S} \;=\; \mathcal{S}_V \times \mathcal{S}_P,
\]
where~\(\mathcal{S}_V\) is the volatile component (kernel data
structures, page cache, in-memory superblock copy, in-memory bitmap)
and~\(\mathcal{S}_P\) is the persistent component (on-disk
superblock, inode table, allocation bitmap, data blocks, RAF journal).
\end{definition}

\begin{definition}[Layout]\label{def:layout}
A \emph{layout}~\(\Lambda\) is a finite map from named regions
(\texttt{sb}, \texttt{inode\_table}, \texttt{bitmap},
\texttt{rs\_journal}, \texttt{data}) to byte ranges in their
hosting layer, together with a redundancy specification (CRC
algorithm, RS parameters) per region.
\end{definition}

The FTRFS v1 layout, formalized in~\cite{desbrieres2026ftrfs}, is one
instance of~\(\Lambda\). BEAMFS targets a strictly extended layout in
which every region is RS-protected and every region of~\(\mathcal{S}_V\)
that mirrors a region of~\(\mathcal{S}_P\) carries the redundancy
metadata necessary to detect drift between the two.

\subsection{Consistency Relation}

\begin{definition}[Per-region consistency]
Let~\(R\) be a region of~\(\Lambda\) with redundancy specification~\(r\).
A state~\(s\) is \emph{R-consistent}, written~\(s \models_R\), iff
applying~\(r\) to the bytes of~\(R\) in~\(s\) yields a zero syndrome
(\textsc{crc} valid and Reed--Solomon syndromes all zero).
\end{definition}

\begin{definition}[Consistency]\label{def:consistent}
A state~\(s = (s_V, s_P) \in \mathcal{S}\) is \emph{consistent},
written~\(s \models\), iff:
\begin{enumerate}
\item For every region~\(R\) of~\(\Lambda_V\) (the volatile-side
  layout), \(s \models_R\).
\item For every region~\(R\) of~\(\Lambda_P\) (the persistent-side
  layout), \(s \models_R\).
\item For every region~\(R\) that has both volatile and persistent
  representations (the superblock, the bitmap, recently accessed
  inodes), the canonical serialization of~\(s_V|_R\) equals
  \(s_P|_R\) byte-for-byte (\emph{coherence}).
\item Every structural sentinel (the zero pads of \texttt{re\_reserved}
  and \texttt{re\_pad} in the v4 RAF journal entry, the validity
  range of \texttt{s\_data\_protection\_scheme}, the high zero
  bytes of feature flags) is at its canonical value.
\end{enumerate}
We write~\(\bot\) for the explicit fail-closed state, distinct from
every consistent state.
\end{definition}

\subsection{The Five Invariants}\label{sub:invariants}

The recovery soundness theorem (\cref{thm:soundness}) requires the
following invariants to hold of the layout~\(\Lambda\) and of the
operator~\(\mathcal{G}\). They are minimal and stated separately so
that an implementation can be audited against them line by line.

\begin{description}
\item[I1 -- Total redundancy.] Every region of~\(\Lambda\) is covered
  by a redundancy specification (CRC for detection, RS for correction)
  whose correction power equals or exceeds the maximum MBU width of
  the targeted technology layer. For a layer with worst-case MBU
  width~\(w\), each RS subblock must tolerate at least~\(w\) symbol
  errors.
\item[I2 -- Sentinel separation.] Each region carries at least one
  byte whose canonical value is fixed and whose corruption is
  detectable in the redundancy syndrome alone, without semantic
  interpretation.
\item[I3 -- Coherence checkpoint.] For any region present in both
  layers, the canonical serialization map is total, deterministic,
  and bijective on its domain of valid values. This rules out
  ambiguous in-memory representations of an on-disk byte sequence.
\item[I4 -- Idempotence.] For every consistent~\(s\),
  \(\mathcal{G}(s) = s\). The recovery operator is the identity on
  consistent states.
\item[I5 -- Fail-closed monotonicity.] If
  \(\mathcal{G}(s) = \bot\) for some~\(s\), then
  \(\mathcal{G}(\bot) = \bot\). Once the operator declares
  fail-closed, no subsequent application reverses that decision.
\end{description}

The first three are properties of the layout; the last two are
properties of the operator. The soundness theorem couples them.

\begin{remark}
A reviewer may legitimately object that I1 is unrealistically
strong: total redundancy increases space overhead, typically by a
factor proportional to~\((n-k)/k\) for an RS\((n, k)\) code (e.g.,
6.7\% for RS(255,239)). We agree, and we treat the total-redundancy
regime as the v1 baseline. A future v2 will explore selective
redundancy under a probabilistic relaxation of the soundness
statement (with bounded silent-corruption risk quantified rather
than zero by construction). The conservative v1 statement is what
allows the present soundness theorem to be \emph{absolute} rather
than statistical.
\end{remark}
